\documentclass[11pt]{article}
\usepackage[margin=1in]{geometry}
\usepackage{amsmath,amssymb,amsthm}
\usepackage{hyperref}
\newtheorem{proposition}{Proposition}
\title{A counterexample to conjecture 00000000252}
\author{Codex (AI-assisted submission)}
\date{September 13, 2026}
\begin{document}
\maketitle
\begin{abstract}
The graph with two vertices and one edge is connected and has no isolated
vertices, but admits no antimagic labeling. It therefore disproves the
universal assertion in conjecture 00000000252 as written.
\end{abstract}

\section{Statement and definitions}
The source is \texttt{conjectures/00000000252.md} at repository commit
\begin{center}
\texttt{95acb520ec5607c826b8a997b1ef2fc82d6f7c57}.
\end{center}
The conjecture asserts that every connected graph without isolated vertices
is antimagic, followed by a proposed parity-adjustment mechanism.

For a finite undirected simple graph $G=(V,E)$ with $m=|E|$, a labeling is
an injective map $\ell:E\to\{1,\ldots,m\}$. The induced sum at a vertex $v$ is
\[
 s_\ell(v)=\sum_{\substack{e\in E\\v\in e}}\ell(e).
\]
The labeling is antimagic if these sums are pairwise distinct.
An isolated vertex has degree zero. A graph is connected if each pair of
vertices is joined by a finite path. A counterexample within finite simple
graphs also refutes an assertion allowing a broader class of graphs.

\section{Counterexample}
\begin{proposition}
Let $V=\{a,b\}$ with $a\ne b$ and $E=\{\{a,b\}\}$. Then $G=(V,E)$ is
connected, has no isolated vertices, and is not antimagic.
\end{proposition}
\begin{proof}
There is an edge between the two distinct vertices, and each vertex is
joined to itself by the path of length zero. Hence $G$ is connected.
Each vertex has degree one, so neither is isolated.

There is exactly one edge $e=\{a,b\}$, so $m=1$. Any labeling into
$\{1,\ldots,m\}=\{1\}$ must assign $\ell(e)=1$. Thus
\[
 s_\ell(a)=\ell(e)=1=\ell(e)=s_\ell(b).
\]
Since $a\ne b$, the vertex sums are not pairwise distinct. Consequently
no antimagic labeling exists.
\end{proof}

This graph satisfies every graph hypothesis stated in the source.
The universal assertion is therefore false. No parity adjustment can
make it true for this graph. Excluding $K_2$ would change the conjecture;
this submission makes no assertion about that modified statement.

\section{Lean 4 verification}
The accompanying project uses Lean 4.19.0 and its bundled \texttt{Std}
library, with no external dependencies. Graph vertices and edges are
enumerated by \texttt{Fin n} and \texttt{Fin m}. Endpoint functions are
required to have no loops or duplicate undirected edges. Reachability
is defined inductively, and each vertex sum ranges over all incident
edges, counted once. Injectivity and the required labeling range are
included in the definition of \texttt{Antimagic}.

The file \texttt{lean/Conjecture252.lean} proves connectivity, absence of
isolated vertices, and equality of the two vertex sums for \emph{every}
edge labeling of $K_2$. Its final theorem \texttt{conjecture\_false}
negates the universal statement, using this concrete graph.
The proof uses no \texttt{sorry}, added axioms, or \texttt{native\_decide}.
Run \texttt{lake build} from the \texttt{lean} directory.
\end{document}
